\documentclass[11pt,a4paper]{article}
\usepackage[margin=23mm]{geometry}
\usepackage[T1]{fontenc}
\usepackage[utf8]{inputenc}
\usepackage{lmodern,microtype,amsmath,amssymb,booktabs,tabularx,array}
\usepackage{graphicx,xcolor,listings,enumitem}
\usepackage[section]{placeins}
\usepackage[colorlinks=true,linkcolor=blue!50!black,urlcolor=blue!50!black]{hyperref}
\usepackage{fancyhdr}
\pagestyle{fancy}
\fancyhf{}
\lhead{Statistical learning: worked notebook explanations}
\rfoot{\thepage}
\setlength{\headheight}{14pt}
\setlength{\parindent}{0pt}
\setlength{\parskip}{6pt}
\setlength{\emergencystretch}{2em}
\setlist{itemsep=3pt,topsep=4pt}
\lstset{language=Python,basicstyle=\small\ttfamily,breaklines=true,
 columns=fullflexible,keepspaces=true,showstringspaces=false,
 frame=single,rulecolor=\color{gray!35},backgroundcolor=\color{gray!4},
 commentstyle=\color{green!35!black},keywordstyle=\color{blue!65!black},
 aboveskip=8pt,belowskip=8pt}
\newcommand{\code}[1]{\texttt{\detokenize{#1}}}
\newcommand{\RSS}{\operatorname{RSS}}
\newcommand{\RMSE}{\operatorname{RMSE}}
\newcommand{\MAE}{\operatorname{MAE}}
\newcommand{\E}{\mathbb{E}}
\newcommand{\ind}[1]{\mathbf{1}\{#1\}}
\newcommand{\fig}[2]{\begin{figure}[htbp]\centering
\includegraphics[width=\linewidth,height=0.39\textheight,keepaspectratio]{figures/#1}
\caption{#2}\end{figure}}
\newcommand{\notebooktitle}[3]{\begin{center}{\LARGE\bfseries #1\par}\vspace{5pt}
{\large Problem statement and worked resolution\par}\vspace{5pt}
Companion to \code{#2}\par\end{center}
\textbf{How to use this explanation.} #3 Section numbers follow the notebook's numbered sections where present. Numerical values are rounded and refer to the recorded run, not universal performance claims.}
\newcommand{\methods}[1]{\subsection*{Main Python calls}#1}

\begin{document}
\notebooktitle{3A: Margins, hinge loss, and support vectors}{03_A_margins_and_hinge.ipynb}{This explanation develops the geometry and optimality checks, interprets the actual results, and distinguishes solver facts from statistical conclusions.}
\section*{Problem and notation}
For binary labels $y\in\{-1,+1\}$, fit a score $f(x)=w^Tx+b$ and predict $+1$ when $f(x)\ge0$.
A functional margin is $yf(x)$, and hinge loss is $\max(0,1-yf(x))$. These are score-based quantities, not probabilities.
We study a five-row exactly solvable example before generating 400 overlapping observations and using cross-validation (CV) to select a linear SVM.
The small example teaches optimization; the generated data teach prediction, where perfect separation is neither promised nor required.

\section{Certify the small solution}
The proposed $w=(1,0)$, $b=0$ has functional margins $(3,1,1,1,3)$ on the five rows.
It satisfies the hard-margin constraints $yf\ge1$. The negative support point requires $w_1-b\ge1$;
the average of two positive support constraints requires $w_1+b\ge1$. Hence $w_1\ge1$, implying
$\|w\|^2/2\ge1/2$. The proposed solution attains the bound, so it is optimal.

Why minimize the norm? Rescaling a score by a positive factor does not move its zero boundary.
Normalizing the closest functional margins to one removes that freedom. The distance between the two
parallel planes $f=-1$ and $f=1$ is then $2/\|w\|$, so minimizing norm maximizes the geometric margin.
Do not call the full width 2 the one-sided margin; the latter is 1 here.

The dual coefficients $(0,1/2,1/4,1/4,0)$ give a second certificate:
\[
\sum_i\alpha_i y_i=0,\quad w=\sum_i\alpha_i y_ix_i,\quad
\alpha_i(y_if_i-1)=0.
\]
Their sum is 1 and their induced vector norm is 1, giving dual value $1-1/2=1/2$.
Primal feasibility, dual feasibility, and matching objective values prove optimality.
The reference \code{SVC} with linear kernel and C=10 matches this solution. It is still a soft-margin solver;
the checks, not the size of C alone, establish the hard-margin result for these rows.
\fig{03_A_margins_and_hinge_1.png}{The zero boundary and the two supporting lines. The three circled training rows have nonzero dual coefficients.}

\section{Generate overlapping classes and separate data roles}
Draw 400 equiprobable labels and independent conditional features $x_1\sim N(1.5y,0.8^2)$,
$x_2\sim N(0,1)$. The first feature carries the class shift, while the second is noise.
With simulation seed 17 and stratified split seed 23, obtain 240 training, 80 validation, and 80 test rows.
Training contains 122 negatives and 118 positives.

\code{StandardScaler} centers and scales coordinates using fitting rows. This determines the geometry
in which the margin penalty is applied. The full-training scaler used in the derivative check is not reused
inside CV: each fold fits its own preprocessing. Otherwise held-out rows would influence the representation.
The marginal distributions overlap, so a correctly implemented SVM can misclassify some future observations.

\section{Understand and check the hinge subgradient}
Using standardized feature rows $z_i$ and $\beta=(b,w_1,w_2)$, the objective is
\[
J(\beta)=\frac1n\sum_i\max(0,1-y_i(w^Tz_i+b))+\frac\lambda2\|w\|^2.
\]
For $m_i=y_i(w^Tz_i+b)<1$, the derivative of the hinge contribution with respect to the score is $-y_i$;
for $m_i>1$ it is zero. Summing active rows gives
\[g_b=-\frac1n\sum_{m_i<1}y_i,\qquad
g_w=\lambda w-\frac1n\sum_{m_i<1}y_i z_i.\]
The denominator is the total number of fitting rows, not the number of active rows. Using the latter would
change the objective and its relative penalty whenever the active set changed. The intercept has no penalty term.

At margin exactly one, the hinge is not differentiable. A row's loss subgradient may be any convex combination
of zero and $-y_i(1,z_i^T)^T$. The implementation chooses zero for that row.
The numerical check deliberately uses a nonzero vector away from every corner and central step $10^{-6}$.
The returned components are approximately $(0.116667,-0.672343,-0.047437)$ and agree with finite differences.
At a corner, disagreement with one chosen subgradient would not automatically indicate an error.

This focused notebook checks the derivative without asking students to tune a second optimization algorithm.
The reference fitting objective is $\|w\|^2/2+C\sum_i\text{hinge}_i$.
Dividing by $nC$ gives the mean-loss convention with $\lambda=1/(nC)$.
For the selected full-training C=1, that is $1/240$. Inside each 192-row CV fit the equivalent lambda is $1/192$,
because this experiment fixes C rather than fixing a mean-loss lambda.

\section{Choose C and interpret the CV comparison}
Each pipeline joins scaling and a linear SVC. \code{cross_val_score} fits fresh copies on four folds and
scores the fifth. Stratification keeps both class proportions approximately consistent.
Balanced accuracy is the average of the negative-class recall and positive-class recall.
\begin{center}\begin{tabular}{rrr}\toprule
C&Mean CV balanced accuracy&Fold SD\\\midrule
0.01&0.970819&0.011414\\0.1&0.962486&0.027159\\1&0.971167&0.018825\\
10&0.958486&0.025688\\100&0.962833&0.027422\\\bottomrule\end{tabular}\end{center}
The rule selects C=1 and refits on all training rows. The tiny difference from C=0.01 is not strong evidence
of an intrinsic advantage. Fold scores share fitting data, and winning a search makes the largest observed
CV value optimistic as a performance summary. Validation and test rows serve separate assessment roles.

\section{Interpret the conflicting observation diagnostic}
Add $(2.5,0)$ with label $-1$ to the training set, keeping the original scaler fixed.
This isolates a fit response to one conflicting example rather than simultaneously changing feature geometry.
The comparison models are temporary; the original selected model remains untouched.
\fig{03_A_margins_and_hinge_2.png}{Solid and dashed boundaries show the original and perturbed fits at two fixed C values. These diagnostic fits are not selected against test data.}
At C=0.01, mean absolute validation score change is about 0.000487. At C=100 it is about 1.13156.
The latter places much more objective weight on hinge loss relative to coefficient norm in this example.
This observation does not prove that increasing C always increases boundary displacement.

Support counts change from 156 to 158 at C=0.01 and from 22 to 29 at C=100.
More support vectors do not necessarily imply worse classification: many can be correctly classified inside
the unit margin. Mean slack also changes, but note that the post-perturbation mean uses 241 rows rather than 240.
Decision-score differences do not directly measure geometric distances unless divided by a norm in a specified coordinate system.

\section{Evaluate the original selected fit}
Validation balanced accuracy is 0.9875. On the final test, the selected model has accuracy and balanced
accuracy 0.9875, ROC AUC 1.0, and confusion counts
\[\begin{pmatrix}40&0\\1&39\end{pmatrix},\]
with true labels as rows and predicted labels as columns. The training-majority baseline gets 0.5 accuracy.
The fitted model uses 33 support vectors, 17 negative and 16 positive.

Perfect AUC alongside one classification error is possible: all positives may rank above all negatives while
the fixed zero cutoff still lies on the wrong side of one positive score. AUC evaluates ordering, whereas
accuracy evaluates one decision rule. The test minimum functional margin is about $-0.5772$, identifying a
wrong-side example; the median is about 2.6941. These margins are not calibrated confidence levels.

\section*{Documentation, interpretation tips, and common errors}
The \href{https://scikit-learn.org/1.8/modules/svm.html}{scikit-learn SVM guide} documents scores, support-vector
attributes, and differences between SVC and LinearSVC. The latter's defaults do not solve the identical
hinge/intercept problem, so substituting it is not an exact solver check.
\begin{itemize}
\item Keep labels in $\{-1,+1\}$ when manually forming $yf$; using 0 for the negative class makes that product meaningless.
\item In binary SVC, inspect \code{classes_} before deciding which class a positive score favors.
\item Support indices refer to fitting rows in their supplied order, not original full-dataset identifiers after arbitrary reshuffling.
\item Scaling affects the model's distance geometry, not just numerical speed. Preserve it with the estimator in a pipeline.
\item A correctly classified point with margin 0.4 has hinge loss 0.6. Classification error and margin violation are different concepts.
\end{itemize}
\end{document}
